\documentclass[a4paper, 14pt]{extarticle}

% Поля
%--------------------------------------
\usepackage{geometry}
\geometry{a4paper,tmargin=2cm,bmargin=2cm,lmargin=3cm,rmargin=1cm}
%--------------------------------------


%Russian-specific packages
%--------------------------------------

\usepackage[T2A]{fontenc}
\usepackage[utf8]{inputenc} 
\usepackage[english, main=russian]{babel}
%--------------------------------------

\usepackage{textcomp}

% Красная строка
%--------------------------------------
\usepackage{indentfirst}               
%--------------------------------------             


%Graphics
%--------------------------------------
\usepackage{graphicx}
\graphicspath{ {./images/} }
\usepackage{wrapfig}
%--------------------------------------

% Полуторный интервал
%--------------------------------------
\linespread{1.3}                    
%--------------------------------------

%Выравнивание и переносы
%--------------------------------------
% Избавляемся от переполнений
\sloppy
% Запрещаем разрыв страницы после первой строки абзаца
\clubpenalty=10000
% Запрещаем разрыв страницы после последней строки абзаца
\widowpenalty=10000
%--------------------------------------

%Списки
\usepackage{enumitem}

%Подписи
\usepackage{caption} 

%Гиперссылки
\usepackage{hyperref}

\hypersetup {
	unicode=true
}

%Рисунки
%--------------------------------------
\DeclareCaptionLabelSeparator*{emdash}{~--- }
\captionsetup[figure]{labelsep=emdash,font=onehalfspacing,position=bottom}
%--------------------------------------


\usepackage{tempora}
\usepackage{amsmath}
\usepackage{color}
\usepackage{listings}
\lstset{
  belowcaptionskip=1\baselineskip,
  breaklines=true,
  frame=L,
  xleftmargin=\parindent,
  language=Python,
  showstringspaces=false,
  basicstyle=\footnotesize\ttfamily,
  keywordstyle=\bfseries\color{blue},
  commentstyle=\itshape\color{purple},
  identifierstyle=\color{black},
  stringstyle=\color{red},
}

%--------------------------------------
%			НАЧАЛО ДОКУМЕНТА
%--------------------------------------

\begin{document}

%--------------------------------------
%			ТИТУЛЬНЫЙ ЛИСТ
%--------------------------------------
\begin{titlepage}
\thispagestyle{empty}
\newpage


%Шапка титульного листа
%--------------------------------------
\vspace*{-60pt}
\hspace{-65pt}
\begin{minipage}{0.3\textwidth}
\hspace*{-20pt}\centering

\includegraphics[width=\textwidth]{emblem}
\end{minipage}
\begin{minipage}{0.67\textwidth}\small \textbf{
\vspace*{-0.7ex}
\hspace*{-6pt}\centerline{Министерство науки и высшего образования Российской Федерации}
\vspace*{-0.7ex}
\centerline{Федеральное государственное бюджетное образовательное учреждение }
\vspace*{-0.7ex}
\centerline{высшего образования}
\vspace*{-0.7ex}
\centerline{<<Московский государственный технический университет}
\vspace*{-0.7ex}
\centerline{имени Н.Э. Баумана}
\vspace*{-0.7ex}
\centerline{(национальный исследовательский университет)>>}
\vspace*{-0.7ex}
\centerline{(МГТУ им. Н.Э. Баумана)}}
\end{minipage}
%--------------------------------------

%Полосы
%--------------------------------------
\vspace{-25pt}
\hspace{-35pt}\rule{\textwidth}{2.3pt}

\vspace*{-20.3pt}
\hspace{-35pt}\rule{\textwidth}{0.4pt}
%--------------------------------------

\vspace{1.5ex}
\hspace{-35pt} \noindent \small ФАКУЛЬТЕТ\hspace{80pt} <<Информатика и системы управления>>

\vspace*{-16pt}
\hspace{47pt}\rule{0.83\textwidth}{0.4pt}

\vspace{0.5ex}
\hspace{-35pt} \noindent \small КАФЕДРА\hspace{50pt} <<Теоретическая информатика и компьютерные технологии>>

\vspace*{-16pt}
\hspace{30pt}\rule{0.866\textwidth}{0.4pt}
  
\vspace{11em}

\begin{center}
\Large {\bf Задача <<Метод прогонки>>} \\ 
\large {\bf по курсу <<Численные методы линейной алгебры>>} \\ 
\end{center}\normalsize

\vspace{8em}


\begin{flushright}
  {Студент группы ИУ9-71Б Лысенко В. В.\hspace*{15pt} \\
  \vspace{2ex}
  Преподаватель Посевин Д. П.\hspace*{15pt}}
\end{flushright}

\bigskip

\vfill
 

\begin{center}
\textsl{Москва 2024}
\end{center}
\end{titlepage}
%--------------------------------------
%		КОНЕЦ ТИТУЛЬНОГО ЛИСТА
%--------------------------------------

\renewcommand{\ttdefault}{pcr}

\setlength{\tabcolsep}{3pt}
\newpage
\setcounter{page}{2}

\section{Задачи}\label{Sect::task}
\par
\renewcommand{\labelenumii}{\arabic{enumi}.\arabic{enumii}.}
\begin{enumerate} 
Реализовать метод прогонки и вычислить погрешность.
\end{enumerate}


\section{Практическая реализация}
\par

Функция progonka() решает систему линейных алгебраических уравнений с трехдиагональной матрицей. Она принимает четыре вектора a, b, c и f как входные параметры, где a, b, c представлюят элементы трёх диагоналей соответственно, а f - результирующий вектор.
\par
Внутри функции progonka() создается два вектора p и q длиной n, которые используются для хранения промежуточных результатов.
\par
Функция norm() вычисляет евклидову норму вектора x, полученного методомм прогонки, а функция err() вычисляет погрешность.

\par
Код программы представлен ниже:
\begin{lstlisting}
n = 1000

function progonka(a, b, c, d)

    p = zeros(n)
    q = zeros(n)

    p[1] = c[1] / b[1]
    q[1] = f[1] / b[1]
    for i in 2:n-1
        k = (b[i] - p[i-1]*a[i])
        p[i] = c[i] / k
        q[i] = (f[i] - q[i-1]*a[i]) / k
    end
    q[n] = (f[n] - q[n-1]*a[n]) / (b[n] - p[n-1]*a[n])

    x = zeros(n)
    x[n] = q[n]
    for i in n-1:-1:1
        x[i] = q[i] - p[i]*x[i+1]
    end
    
    return x
end

function norm(x)
    sum = 0.0
    for element in x
        sum += element^2
    end
    return sqrt(sum)
end

function err(a, b, c, f, x)
    f1 = zeros(n)
    f1[1] = b[1] * x[1] + c[1] * x[2]
    f1[n] = a[n] * x[n-1] + b[n] * x[n]
    for i in 2:n-1
        f1[i] = a[i] * x[i-1] + b[i] * x[i] + c[i] * x[i+1]
    end
    
    for i in 1:n
       # print(f1[i], " ")
    end
    println()
    
    return norm(f1 - f)
end

a = zeros(n)
b = zeros(n)
c = zeros(n)
a[n] = 1.0
c[1] = 1.0
b[1] = 4.0
b[n] = 4.0
for i in 2:n-1
    a[i] = 1.0
    b[i] = 4.0
    c[i] = 1.0
end

f = zeros(n)
#f[1:4] = [5.0, 6.0, 6.0, 5.0]
f[1] = 5.0
f[n] = 5.0
for i in 2:n-1
    f[i] = 6.0
end


x = progonka(a, b, c, f)
for i in 1:n
    #print(x[i], " ")
end
println()

error = err(a, b, c, f, x)
println("Error: ", error)
\end{lstlisting}

\section{Результат}
\begin{center}
\includegraphics[scale=0.75]{g1}
\par
\end{center}

\section{Выводы}
В данной лабораторной работе был реализован метод прогонки и получена погрешность вычислений.
\end{document}
